\documentclass[10pt,a4paper]{amsart}
\usepackage[margin=19mm]{geometry}
\usepackage{amsmath,amssymb,mathtools}
\usepackage[scr=rsfs,cal=euler]{mathalfa}
\usepackage{xfrac}
\usepackage[unicode,hidelinks,bookmarks=false]{hyperref}
\newcommand{\ie}{i.e.\ }
\newcommand{\cf}{cf.\ }
\newcommand{\resp}{respectively\ }
\newcommand{\Subsection}{\subsection}
\providecommand{\nobreakdash}{\nobreak\hbox{-}}
\setlength{\emergencystretch}{2em}
\multlinegap0pt
\usepackage{bm}
\usepackage{bbm}
\usepackage{dsfont}
\usepackage{xspace}
\usepackage{cancel}
\usepackage{mathtools}
\usepackage{amsthm}
\makeatletter
\@ifundefined{lemma}{\newtheorem{lemma}{Lemma}}{}
\@ifundefined{proposition}{\newtheorem{proposition}{Proposition}}{}
\@ifundefined{theorem}{\newtheorem{theorem}{Theorem}}{}
\makeatother
\title[Controllability of periodic linear systems, the Poincare sph]{Controllability of periodic linear systems, the Poincare sphere, and quasi-affine systems}
\author{Fritz Colonius, Alexandre Santana, Juliana Setti}
\date{Source: arXiv:2212.05308v1 (CC BY 4.0)}
\begin{document}
\maketitle

\noindent\emph{Source and licence.} Controllability of periodic linear systems, the Poincare sphere, and quasi-affine systems. Version arXiv:2212.05308v1 (CC BY 4.0), distributed under the Creative Commons Attribution 4.0 International licence, as recorded in the arXiv licence metadata for this exact version and shown on the abstract page https://arxiv.org/abs/2212.05308v1. Full rights notice: licenses/arxiv-2212.05308-CC-BY-4.0.txt.\par\smallskip

\noindent\textit{Editorial scope.} Counted unit: the Theorem whose source label is \texttt{Theorem\_cs1} in the article (source0.tex lines 1222--1241, source label \texttt{Theorem\_cs1}), delivered together with the article's own complete original proof of it (source0.tex lines 1243--1332, one \texttt{proof} environment of 90 lines). No mathematics is written, repaired or re-derived: statement and proof are the article's own text line for line. Required context retained uncounted: periodic0 (lines 70--72); hom (lines 86--88); Proposition9\_JDDE (lines 309--318); proposition\_R (lines 401--428); Theorem7.1.7 (lines 534--565); aug1 (lines 602--605); Theorem\_sub (lines 978--1001); Lemma\_point (lines 1165--1170). Every definition, display and label of those lines is retained unchanged. Omitted from this package and not redistributed: 1--69, 73--85, 89--308, 319--400, 429--533, 566--601, 606--977, 1002--1164, 1171--1221, 1242--1242, 1333--2335. Nothing inside a retained excerpt is omitted or altered: every retained body is delivered exactly as the article's lines are written, so no omission receipt of any kind accompanies this package. Citations: the retained excerpts carry no citation command at all, so no bibliography entry of the article is transcribed and no reference is redistributed. Reference adaptation: none was needed and none was made; every reference inside the delivered unit resolves inside it, so no unresolved reference or citation is produced. Printed slips kept literal: no printed word, symbol, hypothesis or number of the counted statement or of its proof was altered. Layout and class: the article's own class is \texttt{siamltex} with packages including amsfonts, amsmath, amssymb, graphicx, enumerate, color, hyperref; the wrapper is the lane's pinned \texttt{amsart} editorial wrapper at 10pt on A4, which restates the theorem-like environments the retained text needs and adds the article's own macro definitions that the retained text needs (none), with extra wrapper packages (bm, bbm, dsfont, xspace, cancel, mathtools). The delivered numbering is the unit's own. Assets: no retained span contains a figure environment, a graphics-inclusion command, a PDF-page inclusion, a table or a TikZ picture; no image file is copied into this package, the package declares no asset dependency and no image file is redistributed with it. Licence: version arXiv:2212.05308v1 of this article is distributed under the Creative Commons Attribution 4.0 International licence, as recorded in the arXiv licence metadata for this exact version and shown on the abstract page https://arxiv.org/abs/2212.05308v1. A full rights notice is delivered at licenses/arxiv-2212.05308-CC-BY-4.0.txt. Characters: every retained excerpt is pure ASCII, so the delivered engine, XeLaTeX with its default Latin Modern text font, reports no missing character.
\begin{equation}
\dot{x}(t)=A(t)x(t)+B(t)u(t),\quad u(t)\in U, \label{periodic0}%
\end{equation}

\begin{equation}
\dot{x}(t)=A(t)x(t). \label{hom}%
\end{equation}

\begin{proposition}
\label{Proposition9_JDDE}For $\tau\in\left[  0,T\right]  $ consider
$x\in\mathbf{R}_{kT+\tau}(\tau,0)$ and $y\in\mathbf{R}_{\ell T+\tau}(\tau,0)$
where $k,\ell\in\mathbb{N}$. Then it follows that%
\begin{equation}
x+X(kT+\tau,\tau)y=x+X(T+\tau,\tau)^{k}y\in\mathbf{R}_{(k+\ell)T+\tau}%
(\tau,0). \label{2.6}%
\end{equation}

\end{proposition}

\begin{proposition}
\label{proposition_R}Assume that the periodic linear system in
(\ref{periodic0}) without control restrictions is controllable, and let
$\tau\in\left[  0,T\right]  $. Then for system (\ref{periodic0}) with controls
$u\in\mathcal{U}$ the reachable set $\mathbf{R}_{dT+\tau}(\tau,0)$ and the
controllable set $\mathbf{C}_{-dT+\tau}(\tau,0)$ of (\ref{periodic0}) are
convex and contain an $\varepsilon$-ball $\mathbf{B}(0;\varepsilon)$ with
$\varepsilon>0$ around $0\in\mathbb{R}^{d}$. The sets%
\[
\mathbf{R}_{\mathbb{N}T+\tau}(\tau,0):=\bigcup_{k\in\mathbb{N}}\mathbf{R}%
_{kT+\tau}(\tau,0)\text{ and }\mathbf{C}_{-\mathbb{N}T+\tau}(\tau
,0):=\bigcup_{k\in\mathbb{N}}\mathbf{C}_{-kT+\tau}(\tau,0),\tau\in
\lbrack0,T],
\]
are convex and open. Furthermore also the sets%
\[
\mathbf{R}_{\mathbb{N}T+\tau}(0,0),\mathrm{int}\mathbf{R}_{\mathbb{N}T+\tau
}(0,0),\mathbf{C}_{-\mathbb{N}T+\tau}(0,0),\text{ and }\mathrm{int}%
\mathbf{C}_{-\mathbb{N}T+\tau}(0,0)
\]
are convex, and%
\[
\mathbf{R}_{\mathbb{N}T+\tau}(\tau,0)\subset\mathrm{int}\mathbf{R}%
_{\mathbb{N}T+\tau}(0,0)\text{ and }\mathbf{C}_{-\mathbb{N}T+\tau}%
(\tau,0)\subset\mathrm{int}\mathbf{C}_{-\mathbb{N}T+\tau}(0,0).
\]

\end{proposition}

\begin{theorem}
\label{Theorem7.1.7}Let $\Psi=(\theta,\psi):\mathbb{R}\times\mathbb{S}%
^{1}\times\mathbb{R}^{d}\longrightarrow\mathbb{S}^{1}\times\mathbb{R}^{d}$ be
the linear skew product flow associated with the $T$-periodic linear
differential equation (\ref{hom}). For each $\tau\in\mathbb{S}^{1}$ there
exists a decomposition%
\[
\mathbb{R}^{d}=L(\lambda_{1},\tau)\oplus\cdots\oplus L(\lambda_{\ell},\tau)
\]
into linear subspaces $L(\lambda_{j},\tau)$, called the Floquet (or Lyapunov)
spaces, with the following properties:

(i) The Floquet spaces have dimension $d_{j}:=\dim L(\lambda_{j},\tau)$
independent of $\tau\in\mathbb{S}^{1}$.

(ii) They are invariant under multiplication by the principal fundamental
matrix in the following sense:
\[
X(t+\tau,\tau)L(\lambda_{j},\tau)=L(\lambda_{j},\theta(t;\tau))\text{ for all
}t\in\mathbb{R}\text{ and }\tau\in\mathbb{S}^{1}.
\]


(iii) For every $\tau\in\mathbb{S}^{1}$ the Floquet (or Lyapunov) exponents
satisfy%
\[
\lambda(x,\tau):=\lim_{t\rightarrow\pm\infty}\frac{1}{t}\log\Vert\psi
(t;\tau,x)\Vert=\lambda_{j}\text{ if and only if }0\not =x\in L(\lambda
_{j},\tau).
\]

\end{theorem}

\begin{equation}
\dot{\tau}(t)=1\operatorname{mod}T,\quad\dot{x}(t)=A(\tau(t))x(t)+B(\tau
(t))u(t),\qquad u\in\mathcal{U}, \label{aug1}%
\end{equation}

\begin{theorem}
\label{Theorem_sub}Suppose that the periodic system in (\ref{periodic0}) with
unconstrained controls is controllable and consider the autonomized system
(\ref{aug1}) with controls $u\in\mathcal{U}$.

(i) Then the reachable set $\mathbf{R}^{a}(0,0)\subset\mathbb{S}^{1}%
\times\mathbb{R}^{d}$ satisfies $\mathbb{S}^{1}\times\{0\}\subset
\mathrm{int}\mathbf{R}^{a}(0,0)$ and%
\[
X(dT+\cdot,\cdot)\mathcal{K}^{-}\oplus\mathcal{E}^{+,0}\subset\mathrm{int}%
\mathbf{R}^{a}(0,0)\subset\mathcal{K}^{-}\oplus\mathcal{E}^{+,0},
\]
with uniformly bounded convex sets $K_{\tau}^{-}:=\mathrm{int}\mathbf{R}%
_{\mathbb{N}T+\tau}(0,0)\cap E_{\tau}^{-},\tau\in\mathbb{S}^{1}$.

(ii)\ The controllable set $\mathbf{C}^{a}(0,0)$ satisfies $\mathbb{S}%
^{1}\times\{0\}\subset\mathrm{int}\mathbf{C}^{a}(0,0)$ and%
\[
\mathcal{E}^{-,0}\oplus X(-dT+\cdot,\cdot)\mathcal{K}^{+}\subset
\mathrm{int}\mathbf{C}^{a}(0,0)\subset\mathcal{E}^{-,0}\oplus\mathcal{K}^{+}%
\]
with uniformly bounded convex sets $K_{\tau}^{+}:=\mathrm{int}\mathbf{C}%
_{-\mathbb{N}T+\tau}(0,0)\cap E_{\tau}^{+},\tau\in\mathbb{S}^{1}$.
\end{theorem}

\begin{lemma}
\label{Lemma_point}Suppose that the periodic system in (\ref{periodic0}) with
unconstrained controls is controllable. Then $D^{a}:=\overline{\mathbf{R}%
^{a}(0,0)}\cap\mathbf{C}^{a}(0,0)$ is a control set and $\mathbb{S}^{1}%
\times\{0\}\subset\mathrm{int}D^{a}$.
\end{lemma}

\begin{theorem}
\label{Theorem_cs1}Suppose that the periodic system in (\ref{periodic0}) with
unconstrained controls is controllable. Then there exists a unique control set
$D^{a}$ with nonvoid interior of the autonomized system (\ref{aug1}) with
controls $u\in\mathcal{U}$.

It is given by $D^{a}=\overline{\mathbf{R}^{a}(0,0)}\cap\mathbf{C}^{a}(0,0)$
and satisfies $\mathbb{S}^{1}\times\{0\}\subset\mathrm{int}D^{a}$ and, with
$\mathcal{K}^{-}\subset\mathcal{E}^{-}$ and $\mathcal{K}^{+}\subset
\mathcal{E}^{+}$ defined in Theorem \ref{Theorem_sub} and $Y(\tau
):=X(dT+\tau,\tau),\tau\in\mathbb{S}^{1}$,%
\begin{equation}
Y(\cdot)\mathcal{K}^{-}\oplus\mathcal{E}^{0}\oplus Y(\cdot)^{-1}%
\mathcal{K}^{+}\subset\mathrm{int}D^{a}\subset D^{a}\subset Y(\cdot
)^{-1}\overline{\mathcal{K}^{-}}\oplus\mathcal{E}^{0}\oplus Y(\cdot
)\mathcal{K}^{+}. \label{D_a2}%
\end{equation}
In particular, $\mathrm{int}D^{a}$ is unbounded if and only if the center
subbundle $\mathcal{E}^{0}$ is nontrivial.
\end{theorem}

\begin{proof}
The inclusion $\mathbb{S}^{1}\times\{0\}\subset\mathrm{int}D^{a}$ and
$D^{a}=\overline{\mathbf{R}^{a}(0,0)}\cap\mathbf{C}^{a}(0,0)$ follow by Lemma
\ref{Lemma_point}. Furthermore, the inclusions (\ref{D_a2}) imply the last
assertion since $X(dT+\tau,\tau),\tau\in\mathbb{S}^{1}$, as well as
$\mathcal{K}^{-}$ and $\mathcal{K}^{+}$ are bounded. Theorem \ref{Theorem_sub}
implies%
\[
X(dT+\cdot,\cdot)\mathcal{K}^{-}\oplus\mathcal{E}^{+}\oplus\mathcal{E}%
^{0}\subset\mathrm{int}\mathbf{R}^{a}(0,0),\,\mathcal{E}^{-}\oplus
\mathcal{E}^{0}\oplus X(-dT+\cdot,\cdot)\mathcal{K}^{+}\subset\mathrm{int}%
\mathbf{C}^{a}(0,0).
\]
Since $\mathrm{int}D^{a}\subset\mathrm{int}\mathbf{R}^{a}(0,0)\cap
\mathrm{int}\mathbf{C}^{a}(0,0)$ and $X(-dT+\tau,\tau)=X(\tau,-dT+\tau
)^{-1}=X(dT+\tau,\tau)^{-1}$ for $\tau\in\mathbb{S}^{1}$ the first inclusion
in (\ref{D_a2}) follows. In order to prove the third inclusion let
$(\tau,y)\in\mathbf{R}^{a}(0,0)$ be given by%
\[
(\tau,y)=\varphi^{a}(t;(0,0),u)=(t\operatorname{mod}T,\varphi(t;0,0,u)=(\tau
,\varphi(\ell T+\tau;0,u))
\]
with $y=\varphi(\ell T+\tau;0,u)\in\mathbf{R}_{\mathbb{N}T+\tau}(0,0)$. By
Proposition \ref{proposition_R} there is a ball $\mathbf{B}(0;\varepsilon
)\subset\mathbf{R}_{dT+\tau}(\tau,0)$. Hence Proposition
\ref{Proposition9_JDDE} implies%
\[
\mathbf{B}(0;\varepsilon)+X(dT+\tau,\tau)y\subset\mathbf{R}_{(d+\ell)T+\tau
}(\tau,0)\subset\mathbf{R}_{\mathbb{N}T+\tau}(\tau,0).
\]
Since $\varepsilon>0$ is independent of $\tau\in\mathbb{S}^{1}$ it follows
that $X(dT+\cdot,\cdot)\mathbf{R}^{a}(0,0)\subset\mathrm{int}\mathbf{R}%
^{a}(0,0)$, and hence%
\[
X(dT+\cdot,\cdot)\overline{\mathbf{R}^{a}(0,0)}\subset\overline{\mathrm{int}%
\mathbf{R}^{a}(0,0)}.
\]
Analogously it follows that%
\[
X(-dT+\cdot,\cdot)\mathbf{C}^{a}(0,0)\subset\mathrm{int}\mathbf{C}^{a}(0,0).
\]
By Theorem \ref{Theorem_sub} we obtain for $x\in D^{a}=\overline
{\mathbf{R}^{a}(0,0)}\cap\mathbf{C}^{a}(0,0)$,%
\begin{align*}
X(dT+\cdot,\cdot)x  &  \in\overline{\mathrm{int}\mathbf{R}^{a}(0,0)}%
\subset\overline{\mathcal{K}^{-}\oplus\mathcal{E}^{+,0}}=\overline
{\mathcal{K}^{-}}\oplus\mathcal{E}^{0}\oplus\mathcal{E}^{+},\\
X(dT+\cdot,\cdot)^{-1}x  &  =X(-dT+\cdot,\cdot)x\in\mathrm{int}\mathbf{C}%
^{a}(0,0)\subset\mathcal{E}^{-}\oplus\mathcal{E}^{0}\oplus\mathcal{K}^{+}.
\end{align*}
This implies%
\[
x\in X(dT+\cdot,\cdot)^{-1}\left(  \overline{\mathcal{K}^{-}}\oplus
\mathcal{E}^{0}\oplus\mathcal{E}^{+}\right)  \cap X(dT+\cdot,\cdot)\left(
\mathcal{E}^{-}\oplus\mathcal{E}^{0}\oplus\mathcal{K}^{+}\right)  .
\]
By Theorem \ref{Theorem7.1.7} the subbundles $\mathcal{E}^{0}$ and
$\mathcal{E}^{\pm}$ are invariant under $X(dT+\cdot,\cdot)$ and hence%
\[
D^{a}\subset X(dT+\cdot,\cdot)^{-1}\overline{\mathcal{K}^{-}}\oplus
\mathcal{E}^{0}\oplus X(dT+\cdot,\cdot)\mathcal{K}^{+}=Y(\cdot)^{-1}%
\overline{\mathcal{K}^{-}}\oplus\mathcal{E}^{0}\oplus Y(\cdot)\mathcal{K}%
^{+}.
\]
This proves the third inclusion in (\ref{D_a2}).

It remains to show uniqueness. Let $E\subset\mathbb{S}^{1}\times\mathbb{R}%
^{d}$ be an arbitrary control set with nonvoid interior. For all $(\tau,x)\in
E$ there is a control function $u\in\mathcal{U}$ such that $\varphi
^{a}(t;(\tau,x),u)\in E$ for all $t\geq0$, hence $E\cap\left(  \{0\}\times
\mathbb{R}^{d}\right)  \not =\varnothing$.

By linearity, it follows from $\varphi^{a}(t;(0,x_{1}),u)=(\tau,x_{2})$ for
$x_{1},x_{2}\in\mathbb{R}^{d}$ and $t>0$ that $\varphi^{a}(t;\alpha
x_{1},\alpha u)=(\tau,\alpha x_{2})$ for any $\alpha\in(0,1]$. This implies
that $\{(0,\alpha x)\left\vert (0,x)\in E\right.  \}$ is contained in some
control set $D_{\alpha}$ and $\mathrm{int}\{(0,\alpha x)\left\vert (0,x)\in
E\right.  \}\subset\mathrm{int}D_{\alpha}$. Now choose any $(0,x)\in
\mathrm{int}E$ and suppose, by way of contradiction, that
\[
\alpha_{0}:=\inf\{\alpha\in(0,1]\left\vert \forall\beta\in\lbrack
\alpha,1]:(0,\beta x)\in E\right.  \}>0.
\]
Then $(0,\alpha_{0}x)\in\partial E$ and $(0,\alpha_{0}x)\in\mathrm{int}%
D_{\alpha_{0}}$. Therefore $E\cap\mathrm{int}D_{\alpha_{0}}\not =\varnothing$,
and it follows that $E=D_{\alpha_{0}}$ and $(0,\alpha_{0}x)\in\mathrm{int}E$.
This is a contradiction and so $\alpha_{0}=0$. Choosing $\alpha>0$ small
enough such that $(0,\alpha x)\in D^{a}$, we obtain $(0,\alpha x)\in E\cap
D^{a}$ and it follows that $E=D^{a}$.
\end{proof}

\end{document}
